\section{Limit cycles of multiplicity four and their unfolding}
\label{sec:actual-bifurcations}

We continue the fourfold limiting root to a limit cycle of multiplicity four
for small $m>0$ and describe its unfolding together with the two simple fixed
points with $q<0$. We use the standard unfolding in the sign case
$A_1(\lambda_0)=B>0$, $c_0=-C<0$. The proof has three steps: location of all
zeros in $I$, including near the thresholds; construction of the multiple
root; and preparation of the positive cluster in a chart where $D$ is smooth
in $q$.

\begin{lemma}[Values across the thresholds]
\label{mainloc:cross-event}
Set
\begin{equation}\label{mainloc:weights}
 a=\frac{m^4A}{C^3},\qquad b=-\frac{m^2S}{C},\qquad
 \delta=\frac{2m^3T}{3BC^2},\qquad q=\frac{m^2\zeta}{C^2},
 \qquad m>0.
\end{equation}
For every compact set of $(A,S,T)$ and every bounded $\zeta$-interval,
uniformly on their product,
\begin{equation}\label{mainloc:full-limit}
 \frac{C^3}{m^4}D\left(\frac{m^2\zeta}{C^2}\right)
 =A-S\zeta-\zeta^2+\zeta_+^{3/2}
                 +T\sqrt{\zeta_+}+O(\sqrt m).
\end{equation}
This is a $C^0$ estimate, uniform across both thresholds.
\end{lemma}

\begin{proof}
Hold $a,b,m$ fixed and integrate in the separation $t$ from $0$
to $\delta$. Put $y=\Gamma_1(q,t)-t$. The passage
jet and its parameter derivatives give
\[
 \Gamma_{\nu,q}(x)=1+\tfrac32A_\nu\sqrt{x_+}+O(x_+),\qquad
 \Gamma_{\nu,t}(x)=O(x_+^{3/2}),\qquad
 |y-q|\le K(q_+^{3/2}+|t|).
\]
These estimates include zero inputs. Along the integration segment,
$A_2=-B+O(m)$ and $|t|=O(m^3)$. The square-root inequality
$|\sqrt{y_+}-\sqrt{q_+}|\le\sqrt{|y-q|}$ and the total derivative
\eqref{eq:quartic-total-detuning} therefore give
\begin{align*}
 \left|Q_t-\tfrac32B\sqrt{q_+}\right|
 \le K\{&q_+^{3/4}+\sqrt{|t|}+m\sqrt{q_+}+q_++|t|\\
        &+q_+^{3/2}+(q_++|t|)^{3/2}\}=O(m^{3/2}).
\end{align*}
This bound uses $|t|$ and remains valid for $q<0$, so both signs
of separation are included. Integration and the aligned identity
\eqref{eq:passage:aligned-jet} yield
\[
 Q-q=mq_+^{3/2}+\tfrac32B\delta\sqrt{q_+}+O(m^{9/2}),
 \qquad Q-q=O(m^4).
\]
In \eqref{eq:quartic-closing-division}, $c=-C+O(m)$ because the
remaining parameters are fixed at $\xi_0$. Replacing $Q$ by $q$
in $bQ+cQ^2$ changes these terms by $O(m^6)$, while
$(c+C)q^2=O(m^5)$ and $Q^3S=O(m^6)$. Thus
\[
 D=a+bq-Cq^2+mq_+^{3/2}
                    +\tfrac32B\delta\sqrt{q_+}+O(m^{9/2}).
\]
Substitution of \eqref{mainloc:weights} proves the assertion.
\end{proof}

With \eqref{mainloc:weights}, set $\Theta=(A,S,T)$ and
\begin{equation}\label{quartic:centre}
 \Theta_0=\left(-\frac1{256},\frac38,\frac1{16}\right),\qquad
 x_0=\frac14.
\end{equation}
To count cycles on both sides of this bifurcation, we first locate
all zeros near the fourfold configuration.

\begin{lemma}[All roots near the quadruple configuration]
\label{mainloc:all-roots}
There are a neighbourhood $\mathcal O$ of $\Theta_0$, $m_0>0$
and three disjoint compact state intervals $J_{N,-},J_{N,+},J_Q$
about, respectively,
\[
 \zeta_{N,\pm}=\frac{-3\pm2\sqrt2}{16}<0,\qquad
 \zeta_Q=\frac1{16}>0,
\]
such that, for every $\Theta\in\mathcal O$ and $0<m<m_0$, all zeros in
$I$ lie in $(m^2/C^2)(J_{N,-}\cup J_{N,+}\cup J_Q)$.
Each negative interval contains exactly one simple zero. Both passages are
active on the positive interval, which contains all remaining zeros.
The three intervals can be chosen arbitrarily small about their
respective centres, by reducing $\mathcal O$ and $m_0$.
\end{lemma}

\begin{proof}
Bounded parameters satisfy the hypotheses of Lemma~\ref{mainloc:tail}.
Thus all zeros have $|\zeta|\le L$ for a fixed $L$.
At $\Theta_0$, the full-state limit in
\eqref{mainloc:full-limit} is
\[
 -\left(\zeta-\zeta_{N,-}\right)
       \left(\zeta-\zeta_{N,+}\right)\quad(\zeta<0),\qquad
 -\left(\sqrt\zeta-\tfrac14\right)^4\quad(\zeta\ge0).
\]
Its only root positions are the three stated centres, and its value
at zero is $-1/256$. Choose the three small intervals with each centre
in its interval's interior. The complement of their interiors in
$[-L,L]$ is compact and contains no limiting root, so the absolute value
of the limiting function has a positive minimum there. First reducing
$\mathcal O$, then $m_0$, preserves that gap by
Lemma~\ref{mainloc:cross-event}. This excludes all zeros outside the interval
interiors, including their boundaries and the threshold layer.

Take $\mathcal O$ so that $T>0$. For $q<0$ both passages are inactive, so
the closing expansion gives on either negative interval
\[
 \frac{C^3}{m^4}D(m^2\zeta/C^2)
 =A-S\zeta-\zeta^2+O_{C^1}(m).
\]
The limiting derivatives at $\zeta_{N,-},\zeta_{N,+}$ are
$\sqrt2/4,-\sqrt2/4$. Their strict signs and the opposite endpoint
values persist uniformly. The intermediate value theorem and strict
monotonicity give one simple root in each interval. In particular,
their scaled locations are
\[
 \frac{-S\pm\sqrt{S^2+4A}}2+O(m).
\]
Finally the second threshold is
$a_\delta=\Gamma_1^{-1}(\delta)=O(m^3)$, whereas the positive interval
has $q\ge \kappa m^2$ for a fixed $\kappa>0$. Hence that entire interval
lies beyond both thresholds.
\end{proof}

The localization fixes the two negative zeros and places the positive
cluster beyond both thresholds. It remains to meet the four zero conditions
simultaneously and identify the local unfolding. The following theorem refines
Theorem~\ref{thm:actual-quartic-summary}.

\begin{theorem}[Limit cycles of multiplicity four]
\label{thm:actual-quartic}
After a reduction of $\mathcal O$ and $m_0$ in
Lemma~\ref{mainloc:all-roots}, there is a unique smooth branch
$m\mapsto(\Theta_*(m),x_*(m))$, $0\le m<m_0$, equal to $(\Theta_0,x_0)$ at
$m=0$, such that for every $0<m<m_0$ the displacement has a zero of
multiplicity four at $q_*(m)=m^2x_*(m)^2/C^2$ for the parameters
$\Theta_*(m)$. Its parameters and state satisfy
\begin{equation}\label{quartic:physical-branch}
 \begin{aligned}
 a_*(m)&=-\frac{m^4}{256C^3}+O(m^5),&
 b_*(m)&=-\frac{3m^2}{8C}+O(m^3),\\
 \delta_*(m)&=\frac{m^3}{24BC^2}+O(m^4),&
 q_*(m)&=\frac{m^2}{16C^2}+O(m^3).
 \end{aligned}
\end{equation}
Set $p_*(m)=(a_*(m),b_*(m),\xi_0)$. Then
\begin{equation}\label{quartic:curvature-turn-identities}
 q_*(m)=r_{\rm c}(m,p_*(m))^2,\qquad
 \delta_*(m)=\delta_{\rm c}(m,p_*(m)).
\end{equation}
For $0<m<m_0$ there are smooth coordinates $(u,v,w)$ on the parameters
$(a,b,\delta)$, centred on this branch, and a smooth state coordinate $y$,
in which the local zero equation is, up to a nonvanishing smooth factor,
\begin{equation}\label{quartic:preparation}
 \mathfrak p(y)=y^4+uy^2+vy+w=0.
\end{equation}
In the normalized variables $(m,x,A,S,T)$, the changes and the
nonzero factor extend jointly smoothly through $m=0$. Off the multiple-root
set, the number of fixed points in $I$ is two, four or six; each count occurs
on a nonempty open set of parameters. At the fourfold point there are exactly
three fixed points in $I$: the fourfold one and two simple ones with $q<0$.
\end{theorem}

By~\eqref{quartic:curvature-turn-identities}, the fourfold zero lies at the
positive critical point of the separation graph, where $D''=D'''=0$ by
Proposition~\ref{prop:cubic-curvature-window}; the conditions $D=D'=0$ are
additional. The normalized fourfold-root equations construct the branch, and
preparation then gives three independent unfolding parameters on the positive
interval isolated above.

\begin{proof}
Lemma~\ref{lem:weighted-quartic}, after constant rescaling, gives the
jointly smooth auxiliary function
\begin{equation}\label{quartic:joint-quotient}
 F(m,x;\Theta)=\frac{C^3}{m^4}D(m^2x^2/C^2)
 =A-Sx^2+x^3-x^4+Tx+mE_4(m,x,\Theta)
\end{equation}
near $(0,x_0,\Theta_0)$, with all prescribed finite mixed derivatives
bounded. It equals the rescaled displacement for $m>0$.
For $H=(F,F_x,F_{xx},F_{xxx})$, differentiation at that point,
in columns $(A,S,T,x)$, gives
\[
 D_{(A,S,T,x)}H=
 \begin{pmatrix}
 1&-1/16&1/4&0\\
 0&-1/2&1&0\\
 0&-2&0&0\\
 0&0&0&-24
 \end{pmatrix},\qquad \det D_{(A,S,T,x)}H=-48.
\]
The implicit function theorem gives the unique local branch
$\Theta_* =\Theta_0+O(m)$, $x_*=x_0+O(m)$, with
$F_{xxxx}=-24+O(m)\ne0$. Substituting the weights gives
\eqref{quartic:physical-branch}. Lemma~\ref{mainloc:all-roots}
places this zero where both passages are active. For each fixed $m>0$,
$q=m^2x^2/C^2$ is a local diffeomorphism there, so its
multiplicity in $q$ is four.

Put $r_*=\sqrt{q_*}$ and $p_*=(a_*,b_*,\xi_0)$. By
\eqref{quartic:physical-branch},
\[
 r_*=\frac{m}{4C}+O(m^2),\qquad
 \frac{\delta_*}{q_*}=\frac{2m}{3B}+O(m^2).
\]
After reducing $m_0$, this point lies in the outer region of
\eqref{eq:curvature-factor}. At fixed $(m,p,\delta)$, set
$K(r)=r^{-3}\mathcal U(r,m,\delta/r^2,p)\ne0$. That factorization
and differentiation at fixed parameters give
\[
 \begin{aligned}
 D''(r^2)&=K(r)[\delta-\Delta(r,m,p)],\\
 2rD'''(r^2)&=K'(r)[\delta-\Delta(r,m,p)]-K(r)\Delta_r(r,m,p).
 \end{aligned}
\]
At the quadruple point $D''=D'''=0$, hence
$\delta_*=\Delta(r_*,m,p_*)$ and $\Delta_r(r_*,m,p_*)=0$.
Since $r_*>0$ and $c_0(p_*)<0$,
Proposition~\ref{prop:cubic-curvature-window} identifies this critical
point as $r_*=r_{\rm c}(m,p_*)$, proving
\eqref{quartic:curvature-turn-identities}.

The Malgrange preparation theorem~\cite[Chapter~IV, Section~2, Theorem~2.1 and Note~(1), p.~95]{golubitsky1973} in $t=x-x_0$,
with parameters $(m,\Theta)$, followed by a translation removing the cubic term,
gives \eqref{quartic:preparation}. At $m=0$, $\Theta=\Theta_0$, the unit
is $-1$ and the lower coefficients vanish. Writing $\Delta\Theta=\Theta-\Theta_0$,
comparison modulo $t^3$ at $m=0$ gives
\[
 \begin{aligned}
 u&=\Delta S+O(|\Delta\Theta|^2),\\
 v&=2x_0\Delta S-\Delta T+O(|\Delta\Theta|^2),\\
 w&=-\Delta A+x_0^2\Delta S-x_0\Delta T
                                    +O(|\Delta\Theta|^2).
 \end{aligned}
\]
The determinant in $(\Delta A,\Delta S,\Delta T)$ is $1$; equivalently,
$\det D_{(A,S,T)}(F,F_x,F_{xx})=2+O(m)\ne0$ on the branch.
Thus $(m,\Theta)\mapsto(m,u,v,w)$ is a local smooth coordinate change in the
normalized variables through $m=0$, centred on the branch. For each fixed
$m>0$, $(u,v,w)$ are local coordinates on the slice of parameters with this
value of $m$; preparation preserves multiplicities.

Choose the positive interval of Lemma~\ref{mainloc:all-roots} inside this
preparation chart. A common reduction of the parameters confines every real
quartic root there, since the coefficients vanish at the centre. The whole positive
cluster therefore has $0$, $2$ or $4$ simple roots off the multiple-root set.
The polynomials
\[
 y^4+\eta^4,\qquad y^4-\eta^4,\qquad
 (y^2-\eta^2)(y^2-4\eta^2),\qquad 0<\eta\ll1,
\]
realize all three counts on open coefficient sets. The parameter
diffeomorphism and the two simple negative roots give the stated counts,
including three fixed points at the origin.
\end{proof}

\begin{proposition}[Local quartic strata]\label{quartic:strata}
In the coordinates $(u,v,w)$ of \eqref{quartic:preparation},
the multiple-root set is parametrized by
\begin{equation}\label{quartic:double-surface}
 v=-4t^3-2ut,\qquad w=3t^4+ut^2.
\end{equation}
Its triple-root curve, its two-double-root curve and its quadruple
point are respectively
\[
 \begin{aligned}
 (u,v,w)&=(-6t^2,8t^3,-3t^4),&&t\ne0,\\
 (u,v,w)&=(-2t^2,0,t^4),&&t\ne0,\\
 (u,v,w)&=(0,0,0).
 \end{aligned}
\]
The ordinary fold part of \eqref{quartic:double-surface} excludes
both the triple-root curve and a simultaneous second double root.
\end{proposition}

\begin{proof}
Solving $\mathfrak p(t)=\mathfrak p'(t)=0$ gives \eqref{quartic:double-surface}.
On this surface, $\mathfrak p''(t)=0$ gives $\mathfrak p=(y-t)^3(y+3t)$,
and the additional equation $\mathfrak p'''(t)=0$ forces $t=0$.
Two distinct double roots must be $t,-t$, since the cubic coefficient
vanishes, so $\mathfrak p=(y^2-t^2)^2$.
\end{proof}

Figure~\ref{fig:quartic-local-diagram} displays the local root counts in these
coordinates; the coordinate change depends on $m$. These strata describe a
neighbourhood of the fourfold cycle at each fixed small $m>0$.
% Exact discriminant of the prepared quartic, on u=-eps^2.
% Similarity coordinates: y=eps Y, V=v/eps^3, W=w/eps^4.
% Curve parameter: (V,W)=(2t-4t^3,3t^4-t^2).
% Source and independent polynomial checks: scripts/check-local-diagrams.py.
\begin{figure}[!htbp]
\centering
\begin{tikzpicture}[x=6.6cm,y=4.7cm,font=\small,
  boundary/.style={black,line width=.75pt},
  guide/.style={black!55,line width=.35pt},
  leader/.style={black,line width=.4pt,-{Stealth[length=1.4mm]}},
  region/.style={fill=white,inner sep=1.5pt,font=\normalsize}]
  \begin{scope}
    \clip (-.94,-.15) rectangle (.94,.56);
    % The closed central chamber has four simple quartic roots.
    \fill[black!9]
      plot[domain=-.707106781186548:.707106781186548,
        variable=\t,samples=181]
      ({2*\t-4*(\t)^3},{3*(\t)^4-(\t)^2}) -- cycle;
    \draw[boundary]
      plot[domain=-.83:.83,variable=\t,samples=301]
      ({2*\t-4*(\t)^3},{3*(\t)^4-(\t)^2});
  \end{scope}
  \draw[guide,-{Stealth[length=1.5mm]}]
    (-.94,-.15)--(.99,-.15) node[right,black] {$V$};
  \draw[guide,-{Stealth[length=1.5mm]}]
    (-.94,-.15)--(-.94,.60) node[above,black] {$W$};
  \foreach \x/\lab in {-.5/{-\tfrac12},0/0,.5/{\tfrac12}}{
    \draw[guide] (\x,-.15)--(\x,-.164);
    \node[below,font=\scriptsize] at (\x,-.164) {$\lab$};
  }
  \foreach \y/\lab in {0/0,.25/{\tfrac14},.5/{\tfrac12}}{
    \draw[guide] (-.94,\y)--(-.955,\y);
    \node[left,font=\scriptsize] at (-.955,\y) {$\lab$};
  }
  \node[region] at (0,.46) {$2$};
  \node[region,fill=black!9] at (0,.12) {$6$};
  \node[region] at (-.76,.13) {$4$};
  \node[region] at (.76,.13) {$4$};
  \node[region] at (0,-.105) {$4$};
  % Two triple-root cusp points and the two-double-root crossing.
  \draw[boundary,fill=white]
    ({-4/(3*sqrt(6))},{-1/12}) circle[radius=1.6pt];
  \draw[boundary,fill=white]
    ({4/(3*sqrt(6))},{-1/12}) circle[radius=1.6pt];
  \fill (0,.25) circle[radius=1.8pt];
  \node[anchor=east,align=right,font=\scriptsize] (dd)
    at (-.29,.30) {two double\\roots};
  \draw[leader] (dd.east)--(-.015,.252);
  \node[anchor=west,font=\scriptsize] (fold) at (.50,.44) {fold};
  \draw[leader] (fold.west)--(.238,.424);
  \node[anchor=east,font=\scriptsize] (triple) at (-.66,-.055) {triple};
  \draw[leader] (triple.east)--(-.551,-.081);
\end{tikzpicture}
\caption{A slice of the prepared quartic
$y^4+uy^2+vy+w$ at $u=-\varepsilon^2$, for any sufficiently small
$\varepsilon>0$. In the similarity coordinates
$V=v/\varepsilon^3$, $W=w/\varepsilon^4$, the multiple-root set is
$(V,W)=(2t-4t^3,3t^4-t^2)$. Region labels are the total numbers of
simple cycles, including the two simple fixed points with $q<0$.
Open circles mark triple roots; the crossing has two double roots.
This slice does not contain the quadruple-root point $u=v=w=0$.}
\label{fig:quartic-local-diagram}
\end{figure}


\FloatBarrier
